\documentclass[a4paper,11pt]{article}

% Packages
\usepackage[utf8]{inputenc}
\usepackage[T1]{fontenc}
\usepackage{lmodern}
\usepackage{microtype}
\usepackage[french]{babel}
\usepackage{geometry}
\usepackage{xcolor}
\usepackage{enumitem}
\usepackage{titlesec}
\usepackage{tcolorbox}
\usepackage{graphicx}
\usepackage{fancyhdr}
\usepackage{tabularx}
\usepackage{booktabs}
\usepackage{longtable}
\usepackage{array}
\usepackage{tikz}
\usepackage{hyperref}
\usepackage{listings}

% ============================================================
% Couleurs
% ============================================================
\definecolor{mainpurple}{RGB}{106, 90, 205}
\definecolor{lightlavender}{RGB}{230, 230, 250}
\definecolor{deeppurple}{RGB}{75, 0, 130}
\definecolor{softpurple}{RGB}{147, 112, 219}
\definecolor{palelavender}{RGB}{245, 240, 255}
\definecolor{codestring}{RGB}{206, 145, 120}

% ============================================================
% Géométrie
% ============================================================
\geometry{
    top=2.5cm,
    bottom=2.5cm,
    left=2.5cm,
    right=2.5cm,
    headheight=23pt
}

% ============================================================
% En-têtes / pieds de page
% ============================================================
\pagestyle{fancy}
\fancyhf{}
\fancyhead[L]{\color{mainpurple}\small\leftmark}
\fancyhead[R]{\color{softpurple}\small PROJET FTP - SR}
\fancyfoot[C]{\color{mainpurple}\thepage}
\renewcommand{\headrulewidth}{0.5pt}
\renewcommand{\footrulewidth}{0.5pt}
\renewcommand{\headrule}{\hbox to\headwidth{\color{mainpurple}\leaders\hrule height \headrulewidth\hfill}}
\renewcommand{\footrule}{\hbox to\headwidth{\color{mainpurple}\leaders\hrule height \footrulewidth\hfill}}

% ============================================================
% Titres
% ============================================================
\titleformat{\section}
  {\LARGE\bfseries\color{deeppurple}}
  {\thesection}{1em}{}
  [\titlerule]

\titleformat{\subsection}
  {\Large\bfseries\color{mainpurple}}
  {\thesubsection}{1em}{}

\titleformat{\subsubsection}
  {\large\bfseries\color{softpurple}}
  {\thesubsubsection}{1em}{}

% ============================================================
% Boîtes
% ============================================================
\newtcolorbox{infobox}{
    colback=lightlavender,
    colframe=mainpurple,
    boxrule=1pt,
    arc=3mm,
    left=10pt, right=10pt, top=10pt, bottom=10pt
}

\newtcolorbox{tipbox}{
    colback=palelavender,
    colframe=softpurple,
    boxrule=1pt,
    arc=3mm,
    left=10pt, right=10pt, top=10pt, bottom=10pt,
    title={\textbf{Idée}}
}

\newtcolorbox{warningbox}{
    colback=lightlavender!50,
    colframe=deeppurple,
    boxrule=1.5pt,
    arc=3mm,
    left=10pt, right=10pt, top=10pt, bottom=10pt,
    title={\textbf{Attention}}
}

% ============================================================
% Listings
% ============================================================
\lstdefinestyle{codestyle}{
    backgroundcolor=\color{white},
    basicstyle=\ttfamily\small\color{black},
    keywordstyle=\color{mainpurple}\bfseries,
    stringstyle=\color{codestring},
    commentstyle=\color{deeppurple}\itshape,
    breaklines=true,
    breakatwhitespace=false,
    keepspaces=true,
    showspaces=false,
    showstringspaces=false,
    frame=single,
    framesep=4pt,
    rulecolor=\color{mainpurple},
    xleftmargin=8pt,
    xrightmargin=4pt,
    tabsize=4
}

\lstset{style=codestyle}

% ============================================================
% Liens
% ============================================================
\hypersetup{
    colorlinks=true,
    linkcolor=deeppurple,
    citecolor=mainpurple,
    urlcolor=softpurple,
    pdfborder={0 0 0},
    pdftitle={TP FTP -- Compte-rendu},
    pdfauthor={Licence Informatique L3}
}

% ============================================================
% Informations de garde
% ============================================================
\newcommand{\titreprojet}{TP FTP}
\newcommand{\soustitreprojet}{Systèmes et Réseaux}
\newcommand{\nometudiantA}{ANGERETTI Quentin}
\newcommand{\nometudiantB}{CASSOU Baptiste}
\newcommand{\anneeuniversitaire}{2025--2026}

% ============================================================
% Colonnes de tableaux
% ============================================================
\newcolumntype{L}[1]{>{\raggedright\arraybackslash}p{#1}}
\newcolumntype{T}[1]{>{\raggedright\arraybackslash\ttfamily\footnotesize}p{#1}}
\renewcommand{\arraystretch}{1.2}
\setlength{\tabcolsep}{4pt}
\setlength{\LTleft}{0pt}
\setlength{\LTright}{0pt}
\newcommand{\codecell}[1]{{\footnotesize\nolinkurl{#1}}}

\begin{document}

% ============================================================
% PAGE DE TITRE
% ============================================================
\begin{titlepage}
    \centering
    \vspace*{1.2cm}

    {\Huge\bfseries\color{deeppurple} \titreprojet\par}
    \vspace{0.3cm}
    {\LARGE\bfseries\color{mainpurple} \soustitreprojet\par}

    \vspace{1.2cm}

    \begin{tcolorbox}[
        colback=lightlavender,
        colframe=mainpurple,
        width=0.9\textwidth,
        arc=5mm,
        boxrule=2pt
    ]
    \centering
    {\Large\color{deeppurple} Compte rendu du projet de transfert de fichiers\par}
    \vspace{0.35cm}
    {\large Nous décrivons ici le code réel du dépôt, son protocole réseau, son architecture maître / esclaves et les fonctionnalités Q1 à Q17\par}
    \end{tcolorbox}

    \vspace{1.1cm}

    \begin{tikzpicture}
        \fill[mainpurple!30] (0,0) circle (2.8cm);
        \fill[softpurple!40] (0,0) circle (1.9cm);
        \fill[deeppurple!20] (0,0) circle (1.05cm);
        \node[align=center] at (0,0) {\color{deeppurple}\bfseries FTP};
    \end{tikzpicture}

    \vfill

    \begin{tcolorbox}[
        colback=palelavender,
        colframe=softpurple,
        width=0.75\textwidth,
        arc=4mm,
        boxrule=1.2pt
    ]
    \centering
    {\large\textbf{Réalisé par}\par}
    \vspace{0.25cm}
    {\large \nometudiantA\par}
    {\large \nometudiantB\par}
    \vspace{0.35cm}
    {\normalsize Année universitaire : \anneeuniversitaire\par}
    {\normalsize \today\par}
    \end{tcolorbox}
\end{titlepage}

\tableofcontents
\newpage

\section{Introduction}

Nous avons réalisé un petit service de transfert de fichiers inspiré de FTP.
Nous ne cherchons pas à reproduire tout le protocole FTP historique. Nous avons
créé un protocole plus petit, plus simple, et adapté au sujet.

Dans ce projet, nous voulions surtout comprendre quatre idées :

\begin{itemize}[leftmargin=*]
    \item comment un client et un serveur parlent sur TCP ;
    \item comment on définit un protocole applicatif avec des messages fixes ;
    \item comment on transfère un fichier sans accès direct au disque distant ;
    \item comment on ajoute peu à peu des fonctions plus avancées.
\end{itemize}

La version actuelle du dépôt couvre les questions \textbf{Q1 à Q17}. Nous avons
donc :

\begin{itemize}[leftmargin=*]
    \item un client interactif ;
    \item un maître sur le port \texttt{2121} ;
    \item deux serveurs esclaves ;
    \item les commandes \texttt{get}, \texttt{put}, \texttt{ls}, \texttt{rm},
    \texttt{auth} et \texttt{bye} ;
    \item la reprise d'un \texttt{get} interrompu ;
    \item la reconnexion du client en cas de panne pendant un \texttt{get} ;
    \item la réplication simple de \texttt{put} et \texttt{rm} entre esclaves ;
    \item une sérialisation explicite des messages réseau avec prise en compte
    de l'endianness.
\end{itemize}

\begin{infobox}
Ce compte rendu a été réécrit à partir du code actuel. Nous décrivons donc le
fonctionnement réel du dépôt, et pas seulement l'intention du projet.
\end{infobox}
\newpage
\section{Vue d'ensemble}

\subsection{Idée générale}

Le système contient trois programmes :

\begin{itemize}[leftmargin=*]
    \item \texttt{clientFTP} : le programme lancé par l'utilisateur ;
    \item \texttt{masterFTP} : le maître, qui reçoit les clients sur
    \texttt{2121} puis les redirige ;
    \item \texttt{serverFTP} : un esclave, qui traite réellement les commandes
    FTP.
\end{itemize}

Le maître ne transfère pas de fichiers. Il sert seulement à choisir un esclave.
Le vrai travail se fait sur un esclave.

\subsection{Bilan question par question}

\begin{center}
\small
\begin{tabularx}{\textwidth}{@{}L{0.09\textwidth}L{0.16\textwidth}L{0.67\textwidth}@{}}
\toprule
\textbf{Question} & \textbf{État} & \textbf{Ce que nous avons fait} \\
\midrule
Q1 & Réalisée & Nous avons défini le type \texttt{typereq\_t} pour les commandes du protocole. \\
Q2 & Réalisée & Nous avons défini la structure \texttt{request\_t} pour les messages client vers serveur. \\
Q3 & Réalisée & Nous avons créé les binaires client et serveur, ainsi que l'écoute réseau. \\
Q4 & Réalisée & Nous avons géré l'arrêt propre avec \texttt{SIGINT}. \\
Q5 & Réalisée & Nous avons séparé le répertoire client et le répertoire serveur. \\
Q6 & Réalisée & Nous avons reçu une requête, validé la demande et envoyé une réponse. \\
Q7 & Réalisée & Nous avons ajouté la commande interactive \texttt{get} côté client avec statistiques. \\
Q8 & Réalisée & Nous avons envoyé les fichiers par blocs de taille fixe. \\
Q9 & Réalisée & Nous avons gardé une session ouverte pour plusieurs commandes successives. \\
Q10 & Réalisée & Nous avons ajouté la reprise d'un téléchargement et la robustesse après coupure brutale. \\
Q11 & Réalisée & Nous avons fixé un nombre d'esclaves et des ports dédiés. \\
Q12 & Réalisée & Nous avons relié le maître et les esclaves au démarrage. \\
Q13 & Réalisée & Nous avons redirigé le client du maître vers un esclave. \\
Q14 & Réalisée & Nous avons reconnecté automatiquement le client pendant un \texttt{get}. \\
Q15 & Réalisée & Nous avons ajouté la commande \texttt{ls}. \\
Q16 & Réalisée & Nous avons ajouté \texttt{put} et \texttt{rm}, avec réplication simple. \\
Q17 & Réalisée & Nous avons protégé \texttt{put} et \texttt{rm} avec \texttt{auth}. \\
\bottomrule
\end{tabularx}
\end{center}

\subsection{Ce que le code fait aujourd'hui}

La boucle normale est la suivante :

\begin{enumerate}[leftmargin=*]
    \item nous démarrons les esclaves ;
    \item nous démarrons le maître ;
    \item le maître enregistre les esclaves ;
    \item le client se connecte au maître ;
    \item le maître envoie l'adresse d'un esclave ;
    \item le client ouvre une vraie session FTP avec cet esclave ;
    \item le client envoie plusieurs commandes jusqu'à \texttt{bye}.
\end{enumerate}
\newpage
\section{Organisation du dépôt}

\subsection{Binaires principaux}

\begin{center}
\small
\begin{tabularx}{\textwidth}{@{}L{0.20\textwidth}L{0.20\textwidth}L{0.52\textwidth}@{}}
\toprule
\textbf{Binaire} & \textbf{Source principale} & \textbf{Rôle} \\
\midrule
\texttt{bin/clientFTP} & \texttt{src/clientFTP.c} & Lit les commandes utilisateur, se connecte au maître ou à un esclave, puis affiche les résultats. \\
\texttt{bin/masterFTP} & \texttt{src/server\_master.c} & Enregistre les esclaves au démarrage et redirige les clients en round-robin. \\
\texttt{bin/serverFTP} & \texttt{src/serverFTP.c} & Traite les commandes FTP sur un esclave et gère aussi la réplication sur le port de contrôle. \\
\bottomrule
\end{tabularx}
\end{center}

\subsection{Répertoires importants}

\begin{center}
\small
\begin{tabularx}{\textwidth}{@{}L{0.22\textwidth}L{0.70\textwidth}@{}}
\toprule
\textbf{Répertoire} & \textbf{Utilité} \\
\midrule
\texttt{include/} & Tous les fichiers d'en-tête, les constantes, les structures et les prototypes. \\
\texttt{src/} & Le code C du projet. \\
\texttt{data\_client/} & Les fichiers locaux du client. Les téléchargements arrivent ici. Les fichiers envoyés avec \texttt{put} partent de là. \\
\texttt{data\_server/} & Les fichiers vus par les serveurs. Les \texttt{get}, \texttt{ls}, \texttt{put} et \texttt{rm} agissent ici. \\
\texttt{scripts/} & Les scripts de démonstration et de test. \\
\texttt{logs/} & Les logs générés par les scripts. \\
\texttt{bin/} & Les objets compilés et les exécutables. \\
\bottomrule
\end{tabularx}
\end{center}

\begin{warningbox}
Dans ce dépôt, les deux esclaves travaillent dans le même dossier
\texttt{data\_server/}. La séparation entre client et serveur est donc bien
réelle, mais la séparation entre les deux esclaves reste surtout logique.

Autrement dit, nous avons bien un protocole de réplication dans le code, mais
les deux esclaves ne stockent pas leurs fichiers dans deux dossiers physiques
différents.
\end{warningbox}

\subsection{Fichiers les plus importants}

\begin{center}
\small
\begin{tabularx}{\textwidth}{@{}L{0.26\textwidth}L{0.66\textwidth}@{}}
\toprule
\textbf{Fichier} & \textbf{Rôle} \\
\midrule
\texttt{include/ftp\_shared.h} & Définit les constantes, les types de requêtes, les statuts et les structures réseau. \\
\texttt{src/ftp\_protocol.c} & Sérialise, désérialise, valide et envoie les messages du protocole. \\
\texttt{src/ftp\_transfer.c} & Envoie et reçoit les payloads des fichiers et des listes. \\
\texttt{src/clientFTP.c} & Gère l'entrée utilisateur et la connexion côté client. \\
\texttt{src/client\_requests.c} & Implémente \texttt{get}, \texttt{put}, \texttt{ls}, \texttt{rm}, \texttt{auth}, \texttt{bye} côté client. \\
\texttt{src/serverFTP.c} & Démarre un esclave, lance les workers et garde le port de contrôle ouvert. \\
\texttt{src/server\_requests.c} & Implémente les commandes côté serveur et la réplication. \\
\texttt{src/server\_master.c} & Démarre le maître, enregistre les esclaves et répartit les clients. \\
\texttt{src/csapp.c} & Fournit des wrappers utiles pour les sockets, les E/S robustes et les signaux. \\
\texttt{Makefile} & Construit les trois binaires et propose quelques cibles pratiques. \\
\bottomrule
\end{tabularx}
\end{center}
\newpage
\section{Architecture générale}

\subsection{Rôle de chaque programme}

\subsubsection{Le client}

Le client fait trois choses :

\begin{itemize}[leftmargin=*]
    \item il lit une ligne tapée par l'utilisateur ;
    \item il transforme cette ligne en commande ;
    \item il envoie une requête réseau puis affiche le résultat.
\end{itemize}

Le client sait parler soit directement à un esclave, soit au maître. Si nous lui
donnons le port \texttt{2121}, il demande d'abord une redirection au maître.
Sinon, il se connecte directement au port donné.

\subsubsection{Le maître}

Le maître sert de répartiteur de charge très simple.

Il fait deux choses :

\begin{itemize}[leftmargin=*]
    \item au démarrage, il contacte tous les esclaves sur leurs ports de contrôle ;
    \item ensuite, il choisit un esclave pour chaque nouveau client.
\end{itemize}

Le choix se fait en \textbf{round-robin}. Cela veut dire que nous prenons un
esclave, puis le suivant, puis le suivant, et ainsi de suite.

\subsubsection{Les esclaves}

Chaque esclave fait en réalité deux travaux en même temps :

\begin{itemize}[leftmargin=*]
    \item il traite les commandes FTP des clients sur son port client ;
    \item il écoute aussi les messages de réplication sur son port de contrôle.
\end{itemize}

Le père garde le port de contrôle. Les workers traitent les clients.

\subsection{Démarrage du cluster}

Le démarrage suit toujours le même ordre :

\begin{enumerate}[leftmargin=*]
    \item nous lançons un esclave avec un identifiant, par exemple \texttt{1} ;
    \item cet esclave se place dans \texttt{data\_server/} ;
    \item il ouvre son port client et son port de contrôle ;
    \item il attend que le maître se connecte sur le port de contrôle ;
    \item nous faisons la même chose avec le second esclave ;
    \item nous lançons ensuite le maître ;
    \item le maître se connecte aux ports de contrôle \texttt{4001} et
    \texttt{4002} sur \texttt{127.0.0.1} ;
    \item chaque esclave envoie son identité au maître ;
    \item le maître construit la carte du cluster et la renvoie à chaque esclave ;
    \item seulement après cela, le maître écoute sur \texttt{2121}.
\end{enumerate}

Dans la configuration actuelle :

\begin{itemize}[leftmargin=*]
    \item l'esclave 1 écoute pour les clients sur \texttt{3001} et pour le contrôle sur \texttt{4001} ;
    \item l'esclave 2 écoute pour les clients sur \texttt{3002} et pour le contrôle sur \texttt{4002} ;
    \item le maître écoute pour les clients sur \texttt{2121}.
\end{itemize}

\subsection{Démarrage d'une session client}

Quand nous lançons \texttt{clientFTP}, nous faisons d'abord :

\begin{enumerate}[leftmargin=*]
    \item un changement de répertoire vers \texttt{data\_client/} ;
    \item une connexion réseau ;
    \item l'affichage de l'invite \texttt{FTP >>>}.
\end{enumerate}

Ensuite, deux cas existent :

\begin{itemize}[leftmargin=*]
    \item si nous passons un port d'esclave comme \texttt{3001}, le client se connecte directement à cet esclave ;
    \item si nous passons \texttt{2121}, le client se connecte d'abord au maître, reçoit l'adresse d'un esclave, ferme la première connexion, puis ouvre une seconde connexion vers cet esclave.
\end{itemize}
\newpage
\section{Le protocole réseau}

\subsection{Pourquoi nous utilisons TCP}

Nous utilisons TCP parce que nous voulons un flux fiable d'octets.

Pour ce projet, TCP nous apporte :

\begin{itemize}[leftmargin=*]
    \item l'ordre des données ;
    \item la détection des coupures ;
    \item une base simple pour envoyer à la fois des messages de protocole et des fichiers.
\end{itemize}

\begin{tipbox}
TCP n'envoie pas des ``messages'' tout prêts. TCP envoie un flux. Cela veut
dire qu'un \texttt{read()} peut renvoyer moins d'octets que prévu. C'est pour
cela que nous utilisons des fonctions robustes comme \texttt{rio\_readn()} et
\texttt{rio\_writen()}.
\end{tipbox}

\subsection{Messages principaux}

Nous avons défini les messages communs dans \texttt{include/ftp\_shared.h}.

\subsubsection{La requête client}

La structure \texttt{request\_t} contient :

\begin{itemize}[leftmargin=*]
    \item la version du protocole ;
    \item le type de requête ;
    \item un offset pour la reprise ;
    \item une taille de fichier pour \texttt{put} ;
    \item une taille de bloc ;
    \item un nom de fichier ;
    \item un login ;
    \item un mot de passe.
\end{itemize}

\subsubsection{La réponse serveur}

La structure \texttt{response\_t} contient :

\begin{itemize}[leftmargin=*]
    \item un statut ;
    \item le type de requête concerné ;
    \item la taille totale du fichier ;
    \item l'offset accepté par le serveur ;
    \item la taille d'un payload supplémentaire.
\end{itemize}

Le champ \texttt{payload\_size} est utile pour \texttt{ls}, car la réponse
contient alors un texte après l'en-tête.

\subsubsection{Les messages de contrôle}

Nous utilisons aussi d'autres messages :

\begin{itemize}[leftmargin=*]
    \item \texttt{slave\_hello\_t} : identité d'un esclave ;
    \item \texttt{slave\_cluster\_t} : carte complète du cluster ;
    \item \texttt{ftp\_replication\_request\_t} : ordre de réplication pour
    \texttt{put} et \texttt{rm} ;
    \item \texttt{ftp\_control\_reply\_t} : accusé de réception sur le canal de contrôle.
\end{itemize}

\subsection{Ordre des octets et sérialisation}

Un point important du projet est maintenant le suivant :

\begin{warningbox}
Nous n'envoyons plus les structures C brutes sur le réseau.

Nous écrivons chaque champ nous-mêmes dans un buffer réseau. Cela évite les
problèmes de padding mémoire et les problèmes d'endianness entre machines.
\end{warningbox}

Concrètement :

\begin{itemize}[leftmargin=*]
    \item pour les ports réseau, les wrappers sockets utilisent déjà
    \texttt{htons()} ;
    \item pour nos messages applicatifs, nous utilisons \texttt{htonl()} et
    \texttt{ntohl()} pour les champs 32 bits ;
    \item pour les champs 64 bits, nous utilisons un helper dédié du même genre ;
    \item à la réception, nous faisons l'opération inverse avant d'utiliser les valeurs.
\end{itemize}

Nous faisons aussi une validation après décodage :

\begin{itemize}[leftmargin=*]
    \item version du protocole ;
    \item type de requête ;
    \item statut ;
    \item taille de bloc ;
    \item cohérence des offsets et des tailles ;
    \item cohérence des messages de contrôle.
\end{itemize}

Si un message semble avoir été envoyé avec le mauvais ordre des octets, nous le
refusons.

\begin{infobox}
Nous ne convertissons pas le contenu des fichiers eux-mêmes. C'est normal : un
fichier doit voyager comme une suite brute d'octets. Seuls les en-têtes du
protocole doivent être convertis.
\end{infobox}

\subsection{Lecture exacte et écriture exacte}

Le protocole fonctionne parce que nous savons toujours combien d'octets lire :

\begin{itemize}[leftmargin=*]
    \item une requête a une taille fixe sur le réseau ;
    \item une réponse a une taille fixe sur le réseau ;
    \item un message de contrôle a une taille fixe sur le réseau ;
    \item un fichier ou une liste a une taille annoncée dans l'en-tête.
\end{itemize}

Cela simplifie beaucoup le traitement :

\begin{itemize}[leftmargin=*]
    \item nous lisons l'en-tête complet ;
    \item nous interprétons les champs ;
    \item nous lisons ensuite exactement le payload attendu.
\end{itemize}

\subsection{Sécurité minimale sur les noms de fichiers}

La fonction \texttt{ftp\_is\_safe\_filename()} refuse les noms dangereux.

Nous bloquons par exemple :

\begin{itemize}[leftmargin=*]
    \item \texttt{..} ;
    \item \texttt{.} ;
    \item les chemins avec \texttt{/} ;
    \item les noms vides.
\end{itemize}

Cette règle est simple, mais elle évite les sorties du répertoire de travail.
\newpage
\section{Comment une session fonctionne}

\subsection{Boucle générale}

Une session FTP sur un esclave suit cette idée simple :

\begin{enumerate}[leftmargin=*]
    \item le client envoie une requête ;
    \item le serveur la lit ;
    \item le serveur la valide ;
    \item le serveur exécute l'action ;
    \item le serveur envoie une réponse ;
    \item si nécessaire, un payload suit la réponse.
\end{enumerate}

La connexion reste ouverte pour les commandes suivantes. Elle se termine avec
\texttt{bye} ou en cas d'erreur.

\subsection{Commande \texttt{get}}

\subsubsection{Côté client}

Quand nous faisons \texttt{get fichier}, le client :

\begin{enumerate}[leftmargin=*]
    \item vérifie le nom du fichier ;
    \item regarde si un fichier local du même nom existe déjà ;
    \item utilise sa taille locale comme offset de reprise ;
    \item construit une requête \texttt{GET} ;
    \item envoie la requête ;
    \item reçoit la réponse ;
    \item reçoit le reste du fichier et l'écrit localement ;
    \item affiche la taille transférée, le temps et le débit.
\end{enumerate}

\subsubsection{Côté serveur}

Quand un esclave reçoit ce \texttt{GET}, il :

\begin{enumerate}[leftmargin=*]
    \item vérifie la requête ;
    \item fait un \texttt{stat()} sur le fichier demandé ;
    \item ouvre le fichier ;
    \item compare l'offset demandé avec la taille réelle ;
    \item si l'offset est valable, se place au bon endroit ;
    \item envoie une réponse \texttt{OK} ou \texttt{RESTART} ;
    \item envoie ensuite le contenu du fichier par blocs.
\end{enumerate}

\subsubsection{Pourquoi la reprise marche}

La reprise marche parce que :

\begin{itemize}[leftmargin=*]
    \item le client connaît la taille de son fichier local ;
    \item cette taille est envoyée dans \texttt{request\_t.offset} ;
    \item le serveur renvoie l'offset accepté dans \texttt{response\_t.offset} ;
    \item le client écrit ensuite au bon endroit dans le fichier local.
\end{itemize}

\subsection{Commande \texttt{ls}}

Quand nous faisons \texttt{ls}, le serveur :

\begin{enumerate}[leftmargin=*]
    \item parcourt le répertoire courant ;
    \item ignore les noms cachés qui commencent par \texttt{.} ;
    \item construit une grande chaîne avec un nom par ligne ;
    \item envoie une réponse \texttt{OK} ;
    \item envoie ensuite ce texte comme payload.
\end{enumerate}

Le client lit ce texte et l'affiche tel quel.

\subsection{Commande \texttt{put}}

Quand nous faisons \texttt{put fichier}, le client :

\begin{enumerate}[leftmargin=*]
    \item vérifie le nom du fichier ;
    \item fait un \texttt{stat()} local ;
    \item ouvre le fichier à envoyer ;
    \item met sa taille dans la requête ;
    \item envoie la requête ;
    \item attend un premier \texttt{OK} du serveur ;
    \item envoie le contenu du fichier ;
    \item attend un second \texttt{OK}.
\end{enumerate}

Le serveur :

\begin{enumerate}[leftmargin=*]
    \item refuse si la session n'est pas authentifiée ;
    \item envoie un premier \texttt{OK} s'il est prêt ;
    \item reçoit le fichier et l'écrit localement ;
    \item tente de le répliquer vers les autres esclaves ;
    \item envoie enfin un second \texttt{OK} au client.
\end{enumerate}

\subsection{Commande \texttt{rm}}

Quand nous faisons \texttt{rm fichier}, le serveur :

\begin{enumerate}[leftmargin=*]
    \item refuse si la session n'est pas authentifiée ;
    \item supprime le fichier local ;
    \item tente de propager la suppression aux autres esclaves ;
    \item renvoie un statut final.
\end{enumerate}

\subsection{Commande \texttt{auth}}

La commande \texttt{auth <login> <password>} agit sur la session courante.

Le serveur compare les identifiants reçus à deux constantes :

\begin{itemize}[leftmargin=*]
    \item login : \texttt{admin} ;
    \item mot de passe : \texttt{srftp}.
\end{itemize}

Si les identifiants sont bons, la session devient authentifiée. Sinon, elle ne
l'est pas.

\subsection{Commande \texttt{bye}}

La commande \texttt{bye} est simple :

\begin{itemize}[leftmargin=*]
    \item le client envoie une requête \texttt{BYE} ;
    \item le serveur répond \texttt{OK} ;
    \item le client affiche \texttt{Bye!} ;
    \item la connexion se ferme.
\end{itemize}
\newpage
\section{Répartition de charge et réplication}

\subsection{Enregistrement des esclaves}

Chaque esclave envoie au maître un message \texttt{slave\_hello\_t}. Ce message
contient :

\begin{itemize}[leftmargin=*]
    \item la version du protocole ;
    \item l'identifiant de l'esclave ;
    \item son port client ;
    \item son port de contrôle ;
    \item son nom de machine.
\end{itemize}

Le maître vérifie ces informations, puis les garde dans un tableau local.

\subsection{Carte du cluster}

Quand tous les esclaves sont enregistrés, le maître construit un
\texttt{slave\_cluster\_t} qui contient la liste complète des esclaves.

Chaque esclave reçoit cette carte. Grâce à elle, un esclave sait vers quels
autres esclaves il doit envoyer une réplication.

\subsection{Redirection du client}

Quand un client se connecte au maître :

\begin{enumerate}[leftmargin=*]
    \item le maître choisit un esclave ;
    \item il envoie au client le \texttt{slave\_hello\_t} de cet esclave ;
    \item il ferme la connexion ;
    \item le client ouvre ensuite sa vraie session avec cet esclave.
\end{enumerate}

Le choix de l'esclave se fait en round-robin. C'est volontairement simple.

\subsection{Reconnexion après panne}

La reconnexion automatique existe seulement pour \texttt{get}.

Si le client détecte une erreur réseau pendant un \texttt{get}, il :

\begin{enumerate}[leftmargin=*]
    \item ferme la connexion cassée ;
    \item tente de se reconnecter ;
    \item si besoin, repasse par le maître ;
    \item relance le même \texttt{get} ;
    \item reprend le transfert grâce à l'offset local.
\end{enumerate}

Dans la configuration actuelle :

\begin{itemize}[leftmargin=*]
    \item le nombre maximal d'essais vaut \texttt{5} ;
    \item le délai entre deux essais vaut \texttt{1} seconde.
\end{itemize}

\subsection{Réplication de \texttt{put} et \texttt{rm}}

Quand un esclave traite un \texttt{put} ou un \texttt{rm}, il essaie ensuite de
prévenir tous les autres esclaves connus.

\subsubsection{Pour \texttt{put}}

L'esclave source :

\begin{enumerate}[leftmargin=*]
    \item ouvre une connexion vers le port de contrôle du pair ;
    \item envoie une requête de réplication \texttt{FTP\_REPL\_PUT} ;
    \item envoie aussi le contenu du fichier ;
    \item attend un accusé de réception.
\end{enumerate}

L'esclave cible reçoit le fichier, l'écrit localement, puis renvoie un statut.

\subsubsection{Pour \texttt{rm}}

L'esclave source :

\begin{enumerate}[leftmargin=*]
    \item ouvre une connexion vers le pair ;
    \item envoie une requête \texttt{FTP\_REPL\_RM} ;
    \item attend le statut final.
\end{enumerate}

L'esclave cible supprime le fichier puis répond.

\begin{warningbox}
Notre réplication reste simple. Si la modification locale a réussi mais que la
réplication échoue vers un autre esclave, nous écrivons une erreur dans les
logs, mais nous ne revenons pas en arrière.
\end{warningbox}
\newpage
\section{Explication par étapes du sujet}

\subsection{Étape 1 : base du service FTP}

\subsubsection{Q1 et Q2}

Nous avons commencé par définir le langage du protocole :

\begin{itemize}[leftmargin=*]
    \item un type de requête avec \texttt{GET}, \texttt{PUT}, \texttt{LS},
    \texttt{RM}, \texttt{BYE}, \texttt{AUTH} ;
    \item une structure de requête ;
    \item une structure de réponse.
\end{itemize}

Cela a donné une base commune au client et au serveur.

\subsubsection{Q3 et Q4}

Nous avons ensuite mis en place l'exécution :

\begin{itemize}[leftmargin=*]
    \item un client interactif ;
    \item un serveur qui écoute ;
    \item un pool de workers ;
    \item un arrêt propre sur \texttt{SIGINT}.
\end{itemize}

Aujourd'hui, \texttt{NB\_PROC} vaut \texttt{1}. Nous avons donc bien un
mécanisme de workers, mais avec un seul worker par esclave dans la
configuration actuelle.

\subsubsection{Q5}

Nous avons séparé le monde client du monde serveur :

\begin{itemize}[leftmargin=*]
    \item le client travaille dans \texttt{data\_client/} ;
    \item les serveurs travaillent dans \texttt{data\_server/}.
\end{itemize}

Le client ne lit donc pas directement les fichiers du serveur.

\subsubsection{Q6 et Q7}

Nous avons ensuite réalisé le premier vrai transfert :

\begin{itemize}[leftmargin=*]
    \item le client construit un \texttt{GET} ;
    \item le serveur ouvre le fichier demandé ;
    \item le serveur envoie une réponse puis le fichier ;
    \item le client écrit le fichier reçu ;
    \item le client affiche les statistiques.
\end{itemize}

\subsection{Étape 2 : blocs, session longue et reprise}

\subsubsection{Q8}

Nous avons envoyé les fichiers par blocs de taille \texttt{4096} octets. Cela
évite de charger tout le fichier dans un seul gros buffer pour le transfert
normal.

\subsubsection{Q9}

Nous avons gardé une boucle de session. Cela permet d'envoyer plusieurs
commandes sur la même connexion TCP.

Par exemple, nous pouvons faire :

\begin{itemize}[leftmargin=*]
    \item \texttt{get a.txt} ;
    \item \texttt{get b.txt} ;
    \item \texttt{ls} ;
    \item \texttt{bye}.
\end{itemize}

\subsubsection{Q10}

Nous avons traité deux points :

\begin{itemize}[leftmargin=*]
    \item la reprise d'un téléchargement interrompu ;
    \item la robustesse si un client coupe la connexion pendant un envoi.
\end{itemize}

La reprise fonctionne grâce à l'offset. La robustesse fonctionne parce que les
fonctions réseau importantes renvoient des erreurs au lieu de tuer tout le
processus. Le worker peut donc terminer la session courante puis continuer à
servir d'autres clients.

\subsection{Étape 3 : maître et esclaves}

\subsubsection{Q11}

Nous avons fixé les rôles et les ports :

\begin{itemize}[leftmargin=*]
    \item un maître sur \texttt{2121} ;
    \item deux esclaves ;
    \item des ports clients \texttt{3001} et \texttt{3002} ;
    \item des ports de contrôle \texttt{4001} et \texttt{4002}.
\end{itemize}

\subsubsection{Q12}

Nous avons ajouté l'enregistrement des esclaves au démarrage. Le maître apprend
qui sont les esclaves avant d'accepter les clients.

\subsubsection{Q13}

Nous avons ajouté la redirection :

\begin{itemize}[leftmargin=*]
    \item le client parle d'abord au maître ;
    \item le maître choisit un esclave ;
    \item le client parle ensuite à cet esclave.
\end{itemize}

\subsubsection{Q14}

Nous avons ajouté la reconnexion automatique pendant \texttt{get}. Si une
connexion tombe, le client peut :

\begin{itemize}[leftmargin=*]
    \item se reconnecter ;
    \item recevoir une nouvelle redirection ;
    \item reprendre le téléchargement.
\end{itemize}

\subsection{Étape 4 : commandes avancées}

\subsubsection{Q15}

Nous avons ajouté \texttt{ls}. Le serveur envoie une liste texte des fichiers du
répertoire courant.

\subsubsection{Q16}

Nous avons ajouté :

\begin{itemize}[leftmargin=*]
    \item \texttt{put} pour envoyer un fichier du client vers le serveur ;
    \item \texttt{rm} pour supprimer un fichier côté serveur ;
    \item la réplication simple de ces modifications vers les autres esclaves.
\end{itemize}

\subsubsection{Q17}

Nous avons ajouté une authentification de session. Sans authentification :

\begin{itemize}[leftmargin=*]
    \item \texttt{get} et \texttt{ls} restent autorisés ;
    \item \texttt{put} et \texttt{rm} sont refusés.
\end{itemize}
\newpage
\section{Fonctions importantes}

\subsection{Fonctions côté client}

\begin{center}
\small
\begin{tabularx}{\textwidth}{@{}L{0.26\textwidth}L{0.22\textwidth}L{0.44\textwidth}@{}}
\toprule
\textbf{Fonction} & \textbf{Fichier} & \textbf{Rôle simple} \\
\midrule
\codecell{parse_command} & \codecell{src/clientFTP.c} & Analyse la ligne tapée par l'utilisateur. \\
\codecell{ftp_client_connect} & \codecell{src/clientFTP.c} & Se connecte soit au maître, soit directement à un esclave. \\
\codecell{ftp_client_reconnect} & \codecell{src/clientFTP.c} & Réessaie un \texttt{get} après une panne réseau. \\
\codecell{ftp_client_run} & \codecell{src/clientFTP.c} & Gère toute la boucle interactive du client. \\
\codecell{ftp_client_get} & \codecell{src/client_requests.c} & Envoie un \texttt{GET}, reçoit la réponse, reçoit le fichier et calcule les statistiques. \\
\codecell{ftp_client_put} & \codecell{src/client_requests.c} & Envoie un \texttt{PUT} en deux temps. \\
\codecell{ftp_client_ls} & \codecell{src/client_requests.c} & Envoie \texttt{LS} et reçoit le texte de la liste. \\
\codecell{ftp_client_rm} & \codecell{src/client_requests.c} & Envoie \texttt{RM} et lit le statut final. \\
\codecell{ftp_client_auth} & \codecell{src/client_requests.c} & Envoie les identifiants pour authentifier la session. \\
\codecell{ftp_client_bye} & \codecell{src/client_requests.c} & Termine proprement la session. \\
\bottomrule
\end{tabularx}
\end{center}

\subsection{Fonctions côté serveur}

\begin{center}
\small
\begin{tabularx}{\textwidth}{@{}L{0.28\textwidth}L{0.22\textwidth}L{0.42\textwidth}@{}}
\toprule
\textbf{Fonction} & \textbf{Fichier} & \textbf{Rôle simple} \\
\midrule
\codecell{ftp_server_run} & \codecell{src/serverFTP.c} & Démarre un esclave, ouvre les ports et lance les workers. \\
\codecell{register_slave_to_master} & \codecell{src/serverFTP.c} & Attend le maître puis envoie l'identité de l'esclave. \\
\codecell{worker_loop} & \codecell{src/serverFTP.c} & Attend un client et traite sa session. \\
\codecell{control_loop} & \codecell{src/serverFTP.c} & Attend les messages de réplication. \\
\codecell{ftp_handle_client} & \codecell{src/server_requests.c} & Boucle principale d'une session côté serveur. \\
\codecell{handle_get_request} & \codecell{src/server_requests.c} & Gère le téléchargement avec reprise. \\
\codecell{handle_put_request} & \codecell{src/server_requests.c} & Gère le téléversement et déclenche la réplication. \\
\codecell{handle_ls_request} & \codecell{src/server_requests.c} & Gère la liste des fichiers. \\
\codecell{handle_rm_request} & \codecell{src/server_requests.c} & Supprime un fichier puis tente de répliquer la suppression. \\
\codecell{handle_auth_request} & \codecell{src/server_requests.c} & Valide ou refuse l'authentification. \\
\codecell{ftp_handle_replication_connection} & \codecell{src/server_requests.c} & Traite un message de réplication reçu depuis un autre esclave. \\
\bottomrule
\end{tabularx}
\end{center}

\subsection{Fonctions du protocole et du transfert}

\begin{center}
\small
\begin{tabularx}{\textwidth}{@{}L{0.30\textwidth}L{0.22\textwidth}L{0.40\textwidth}@{}}
\toprule
\textbf{Fonction} & \textbf{Fichier} & \textbf{Rôle simple} \\
\midrule
\codecell{ftp_send_request} / \codecell{ftp_receive_request} & \codecell{src/ftp_protocol.c} & Envoient et lisent une requête complète au bon format réseau. \\
\codecell{ftp_send_response} / \codecell{ftp_receive_response} & \codecell{src/ftp_protocol.c} & Envoient et lisent une réponse complète au bon format réseau. \\
\codecell{ftp_send_slave_hello} / \codecell{ftp_receive_slave_hello} & \codecell{src/ftp_protocol.c} & Gèrent l'échange d'identité entre maître et esclaves. \\
\codecell{ftp_send_cluster_map} / \codecell{ftp_receive_cluster_map} & \codecell{src/ftp_protocol.c} & Gèrent l'envoi de la carte du cluster. \\
\codecell{ftp_send_replication_request} / \codecell{ftp_receive_replication_request} & \codecell{src/ftp_protocol.c} & Gèrent les messages de réplication. \\
\codecell{ftp_send_control_reply} / \codecell{ftp_receive_control_reply} & \codecell{src/ftp_protocol.c} & Gèrent les accusés de réception du canal de contrôle. \\
\codecell{ftp_send_fd_payload} & \codecell{src/ftp_transfer.c} & Envoie le contenu d'un fichier ouvert par blocs. \\
\codecell{ftp_receive_file_payload} & \codecell{src/ftp_transfer.c} & Reçoit un fichier et l'écrit sur disque, avec reprise si besoin. \\
\codecell{ftp_receive_buffer_payload} & \codecell{src/ftp_transfer.c} & Reçoit un texte, par exemple le résultat de \texttt{ls}. \\
\codecell{ftp_is_safe_filename} & \codecell{src/ftp_protocol.c} & Refuse les noms de fichiers dangereux. \\
\bottomrule
\end{tabularx}
\end{center}

\subsection{Wrappers système utiles}

Nous utilisons surtout quelques wrappers de \texttt{csapp} :

\begin{itemize}[leftmargin=*]
    \item \texttt{Open\_listenfd} pour ouvrir une socket d'écoute ;
    \item \texttt{open\_clientfd} pour ouvrir une connexion TCP ;
    \item \texttt{rio\_readn} et \texttt{rio\_writen} pour lire et écrire exactement le bon nombre d'octets ;
    \item \texttt{Fork} pour créer les workers ;
    \item \texttt{Signal} pour installer les handlers ;
    \item \texttt{Close} pour fermer proprement les descripteurs.
\end{itemize}
\newpage
\section{Tests}

\subsection{Scripts présents dans le dépôt}

Le dépôt contient plusieurs scripts de test et de démonstration. Les deux plus
importants pour la version finale sont :

\begin{itemize}[leftmargin=*]
    \item \texttt{scripts/demo\_q1\_q17.sh} ;
    \item \texttt{scripts/test\_questions\_q1\_q17.sh}.
\end{itemize}

\subsection{Ce que ces scripts vérifient}

Ces scripts vérifient notamment :

\begin{itemize}[leftmargin=*]
    \item la compilation du projet ;
    \item le démarrage du cluster maître / esclaves ;
    \item un \texttt{get} simple ;
    \item plusieurs commandes sur une seule session ;
    \item la reprise d'un fichier tronqué ;
    \item la robustesse après coupure brutale d'un client ;
    \item la redirection par le maître ;
    \item la reconnexion du client après panne d'un worker ;
    \item \texttt{ls}, \texttt{auth}, \texttt{put} et \texttt{rm} ;
    \item les traces de réplication dans les logs.
\end{itemize}

\subsection{Méthodes de vérification}

Les scripts utilisent des outils simples :

\begin{itemize}[leftmargin=*]
    \item \texttt{grep} pour chercher des messages dans les logs ;
    \item \texttt{test -f} pour vérifier qu'un fichier existe ;
    \item \texttt{sha256sum} pour comparer deux fichiers ;
    \item des attentes explicites dans les logs pour vérifier que le cluster est prêt.
\end{itemize}

\subsection{Commandes utiles}

Les commandes les plus utiles sont :

\begin{itemize}[leftmargin=*]
    \item \texttt{make all} pour construire le projet ;
    \item \texttt{make clean} pour nettoyer ;
    \item \texttt{make demo-q1-q17} pour lancer la démonstration complète ;
    \item \texttt{bash scripts/test\_questions\_q1\_q17.sh} pour lancer la batterie de tests question par question.
\end{itemize}
\newpage
\section{Limites actuelles}

Le projet répond au sujet, mais nous avons gardé quelques choix simples :

\begin{itemize}[leftmargin=*]
    \item les identifiants d'authentification sont codés en dur ;
    \item l'authentification vaut seulement pour la connexion TCP courante ;
    \item la reconnexion automatique existe seulement pour \texttt{get} ;
    \item le maître fait un round-robin simple, sans mesure réelle de charge ;
    \item la réplication est simple et sans journal persistant ;
    \item si un pair est absent pendant une réplication, nous n'avons pas de rattrapage différé ;
    \item les deux esclaves partagent le même dossier \texttt{data\_server/} dans ce dépôt local ;
    \item la taille du pool vaut \texttt{1} par défaut, même si la structure du code accepte plusieurs workers.
\end{itemize}

\newpage
\section{Conclusion}

Nous avons construit un mini service FTP complet pour le sujet.

Le projet montre plusieurs points importants :

\begin{itemize}[leftmargin=*]
    \item un protocole applicatif simple sur TCP ;
    \item des transferts de fichiers par blocs ;
    \item la reprise d'un téléchargement ;
    \item une architecture maître / esclaves ;
    \item une redirection client vers un esclave ;
    \item une reconnexion côté client ;
    \item des commandes avancées avec authentification et réplication simple ;
    \item une sérialisation explicite des messages réseau avec gestion de l'ordre des octets.
\end{itemize}

Nous avons volontairement gardé un design lisible. Ce n'est pas un FTP complet
au sens industriel, mais c'est un projet cohérent, progressif et utile pour
comprendre les bases des systèmes répartis et des réseaux.

\end{document}
